\documentclass[sigconf,anonymous=false]{acmart}

\usepackage{booktabs}
\usepackage{multirow}
\usepackage{amsmath}
% amssymb removed: newtxmath (loaded by acmart) provides blackboard bold, avoids \Bbbk conflict
\usepackage{graphicx}
\usepackage{algorithm}
\usepackage{algpseudocode}
\usepackage{xcolor}
\usepackage{listings}
\usepackage{hyperref}

\title{A Measurement Audit of a Quantum-Classical Feedback Loop on Superconducting QPUs: \\ The Quantum-Hydrodynamic Feedback (QHF) Framework}

\author{Máté Róbert}
\authornote{Corresponding author. Open to research collaboration: \texttt{materobi86@gmail.com}}
\orcid{0009-0003-3986-6039}
\email{materobi86@gmail.com}
\affiliation{
  \institution{Independent Researcher}
  \streetaddress{Budapest}
  \country{Hungary}
}

\author{Szilvia Steiner}
\affiliation{
  \institution{Independent Researcher}
  \country{Hungary}
}

\renewcommand{\shortauthors}{Róbert and Steiner}

\begin{abstract}
We present a \emph{measurement audit} of the Quantum-Hydrodynamic Feedback (QHF) framework: a quantum-classical feedback system coupling a Gross-Pitaevskii fluid solver (a 16-mode observable space) with a 2-qubit quantum bridge on IBM 156-qubit superconducting QPUs. The system implements three separate quantum measurements ($A$, $f$, $\varphi$) per iteration, with real-time Bell-parity anchoring and closed-loop decoding via a Rust bridge. We report the full measurement record --- including the negative results --- and we state explicitly which claims the hardware supports and which it does not. We report 111 documented hardware iterations (68+20+20+3) on \texttt{ibm\_fez} under the IBM Quantum open plan (600 s quota, no Runtime Sessions), measuring queue wait (17.9 s/iter), QPU time (2.96 s/iter), mode distributions, and per-channel decoding accuracy. Four findings are reported as measurements rather than claims: the quantum path achieves 0/6 mode-argmax accuracy (Wilson 95\% CI [0, 0.39]) against a dominant classical reference; the $c$-indicator (Bell balance) tracks hardware drift but \emph{fails} as a decoding correction for amplitude, because the readout noise is asymmetric and the symmetric $p = c\,p + (1-c)/2$ model does not hold; the 16-mode energy gap (1.85\%) is below the shot-noise floor (0.137$\sigma$ at 1024 shots), so the mode competition is not resolvable on current hardware; and the mode histograms of the 2026-10-27 runs are \emph{not} evidence for the quantum channel --- replayed through the corrected fluid at the corrected frequency scale, the hardware's own recorded measurements land inside the band that measurement noise alone produces on all four statistics, so the multi-mode pattern of those runs was produced by a protocol defect rather than by the quantum bridge. Complementary validation: Quantum Delayed Synchronization (QDS) Hahn-echo $T_2 = 89.63 \pm 39.18\ \mu$s (5 runs, SD over runs); Quantum Anchor protocol for coherent-state anchoring; Wave-Quantum frequency-to-correlation encoding (Pearson 0.87--0.999); Quantum Anchored Memory (QAM) integrating $\psi(t)$ dynamics with a 9-level zoom microscope memory (Merkle integrity, spine bus). We release the raw data (126 files, of which 111 are iteration records), analysis scripts, and open-source crates (\texttt{fluid-stratos}, \texttt{quantum-bridge}) under MIT license.
\end{abstract}

\begin{CCSXML}
<ccs2012>
<concept>
<concept_id>10010520.10010553.10010562</concept_id>
<concept_desc>Computer systems organization~Quantum computing</concept_desc>
<concept_significance>500</concept_significance>
</concept>
<concept>
<concept_id>10010147.10010178.10010205</concept_id>
<concept_desc>Computing methodologies~Cognitive architectures</concept_desc>
<concept_significance>500</concept_significance>
</concept>
<concept>
<concept_id>10010520.10010553.10010554</concept_id>
<concept_desc>Computer systems organization~Quantum error mitigation</concept_desc>
<concept_significance>300</concept_significance>
</concept>
<concept>
<concept_id>10010520.10010553.10010555</concept_id>
<concept_desc>Computer systems organization~Quantum algorithms</concept_desc>
<concept_significance>300</concept_significance>
</concept>
</ccs2012>
\end{CCSXML}

\ccsdesc[500]{Computer systems organization~Quantum computing}
\ccsdesc[500]{Computing methodologies~Cognitive architectures}
\ccsdesc[300]{Computer systems organization~Quantum error mitigation}
\ccsdesc[300]{Computer systems organization~Quantum algorithms}

\keywords{Quantum-classical feedback, Gross-Pitaevskii equation, superconducting QPU, hardware validation, neuromorphic memory, quantum error mitigation}

\begin{document}

\maketitle

\section{Introduction}

Quantum-classical hybrid algorithms promise advantage by delegating structured subroutines to quantum hardware while retaining classical control. Yet most demonstrations remain in simulation or use shallow circuits without sustained feedback. We ask: \emph{can a quantum-classical loop run on real hardware long enough to exhibit dynamical structure, and what does the hardware teach us about the loop's viability?}

The Quantum-Hydrodynamic Feedback (QHF) framework answers this by integrating:
\begin{enumerate}
\item \textbf{Fluid STRATOS}: a 2D Gross-Pitaevskii equation (GPE) solver with real split-step Fourier kinetics (\texttt{rustfft 6.4.1}, Strang splitting), hosting 16 cognitive modes as attractors.
\item \textbf{Quantum Bridge}: 2 qubits, 3 separate measurements per iteration --- amplitude ($A$: RY), frequency ($f$: RY+CNOT), phase ($\varphi$: RY+CNOT) --- decoded via analytical inverses.
\item \textbf{Hardware Loop}: circuit arraying (3 circuits + Bell anchor per job), live fidelity tracking ($c$-indicator), quota-aware submission, timeout handling.
\end{enumerate}

We augment this core with independently validated modules:
\begin{itemize}
\item \textbf{QDS} (Quantum Delayed Synchronization): Hahn-echo $T_2$ sweep on \texttt{ibm\_kingston} qubit 91, 5 runs, raw counts.
\item \textbf{Quantum Anchor}: post-pulse protocol for coherent-state anchoring on IBM/IQM.
\item \textbf{Wave-Quantum}: frequency-to-correlation encoding, X-basis ZZ theorem, hardware-validated no-go.
\item \textbf{QAM} (Quantum Anchored Memory): $\psi(t) = A e^{-\gamma t} \cos(2\pi f t + \varphi)$ dynamics fused with Microscope Memory (9-level zoom D0--D8, 13-layer consciousness, mmap, SIMD, spine bridge REST/MCP).
\item \textbf{scs-quantum}: TKR frequency-domain wave-packet encoding, Bell balance 99.3\%.
\end{itemize}

This paper reports the \emph{complete measured record}: 111 QPU iterations on \texttt{ibm\_fez}, QDS $T_2$ data, and the architectural synthesis into QAM. All raw data, scripts, and crates are version-controlled and reproducible.

\section{Related Work}

\paragraph{Quantum-classical loops.} Variational algorithms (VQE, QAOA) iterate parameter updates classically \cite{Peruzzo2014, Farhi2014}. Quantum kernels embed data in Hilbert space \cite{Havlicek2019}. Our loop differs: the quantum subroutine \emph{measures} a wavepacket $(A,f,\varphi)$ injected by classical dynamics, and the decoded values \emph{perturb} the fluid --- a dynamical coupling, not optimization.

\paragraph{Gross-Pitaevskii in cognition.} GPE models appear in Bose-Einstein condensates \cite{Pitaevskii2003} and neural field theories \cite{Coombes2005}. Fluid STRATOS implements 16 modes on a 2D lattice with density redistribution (max $|\Delta| = 2.780$, coherence $0.452$ after 500 steps \cite{QOI-gpe}).

\paragraph{Hardware-validated quantum memory.} Quantum RAM proposals \cite{Giovannetti2008} remain theoretical. QAM grounds $\psi(t)$ in measured $T_2$ and anchor protocols, with classical Microscope Memory providing durable, queryable storage.

\paragraph{Error mitigation.} Zero-noise extrapolation \cite{Temme2017}, probabilistic error cancellation \cite{Endo2018}, and measurement calibration \cite{Nation2019} address symmetric noise. Our $c$-indicator exposes \emph{asymmetric} readout error (floor/ceiling at $P\approx 0.037/0.93$) and bimodal $\varphi$ dynamics, limiting correction efficacy.

\section{System Architecture}

\subsection{Fluid STRATOS (GPE Solver)}

The 2D GPE on $N\times N$ grid ($N=128$ default in the Rust crate, $N=64$ in the FFT validation script, $N=32$ in the loop simulation):
\begin{equation}
i\frac{\partial\psi}{\partial t} = -\frac{1}{2}\nabla^2\psi + g|\psi|^2\psi + V\psi
\end{equation}
Strang splitting: $K/2 \to V \to K/2$ with cached FFT plans. Kinetic step in Fourier space:
\begin{equation}
\psi(\mathbf{k}, t+\Delta t/2) = \psi(\mathbf{k}, t) \exp\left(-i \frac{|\mathbf{k}|^2}{2} \frac{\Delta t}{2}\right)
\end{equation}
Thread-local scratch buffers (\texttt{RefCell<Vec<Complex<f64>>>}) avoid allocation in parallel region. 30 regression tests pass (12 fluid-stratos + 18 quantum-bridge); clippy clean.

\subsection{Quantum Bridge Encoding}

The bridge and its decoding mathematics live in one place, the Rust crate
\texttt{crates/quantum-bridge} \cite{QOI-bridge}, so that the loop, the standalone
measurement script and the simulation cannot drift apart.

Three independent circuits per iteration (2 qubits, 1024 shots each):

\begin{table}[h]
\centering
\caption{Encoding circuits. $\theta_A = A$, $\theta_f = 2\pi f$, $\theta_\varphi = \varphi$.}
\label{tab:encoding}
\begin{tabular}{lccc}
\toprule
Channel & Qubits & Circuit & Ideal $P(q_0=1)$ \\
\midrule
$A$ (amplitude) & 1 & $\text{RY}(\theta_A)$ & $\sin^2(A/2)$ \\
$f$ (frequency) & 2 & $\text{RY}(\theta_f) \to \text{CNOT}$ & $\sin^2(\pi f)$ \\
$\varphi$ (phase) & 2 & $\text{RY}(\theta_\varphi) \to \text{CNOT}$ & $\sin^2(\varphi/2)$ \\
\bottomrule
\end{tabular}
\end{table}

Decoding (analytical, implemented in Rust \texttt{bridge-cli}):
\begin{align}
A &= 2 \arccos(\sqrt{P(q_0=0)}) \\
\varphi &= 2 \arcsin(\sqrt{P(q_0=1)}) \\
f &= \arcsin(\sqrt{P(q_0=1)}) / \pi
\end{align}

No RZ gates (virtual Z not calibrated on open plan). Degradation model for simulation:
\begin{equation}
p_\text{meas} = c \cdot p_\text{ideal} + (1-c) \cdot 0.5
\end{equation}
with inverse \texttt{unbias\_probability} (clamped to $[0,1]$). 18/18 tests pass.

\subsection{Hardware Loop Protocol}

\begin{algorithm}
\caption{Quantum-Classical Hardware Loop}
\begin{algorithmic}[1]
\State $\psi_0 \gets \text{FluidStratos}.\text{init}()$
\For{$i = 1 \to N_\text{iter}$}
    \State Extract $(A_i, f_i, \varphi_i)$ from dominant mode of $\psi_{i-1}$
    \State Submit job: 3 circuits $(A_i, f_i, \varphi_i)$ + Bell anchor (if $i \bmod 5 = 0$)
    \State Wait for job completion (timeout 900 s)
    \State Decode $(A_i', f_i', \varphi_i')$ via \texttt{bridge-cli} using live $c$
    \State $\psi_i \gets \text{FluidStratos}.\text{tick}(\psi_{i-1}, A_i', f_i', \varphi_i')$
    \If{$i \bmod 5 = 0$} \State Update $c \gets \text{Bell parity from anchor}$ \EndIf
\EndFor
\end{algorithmic}
\end{algorithm}

Circuit arraying bundles 4 circuits/job (3 signal + 1 Bell), reducing queue overhead. Pre-submission quota check via \texttt{service.usage()} avoids silent queue-stuck (exhausted quota accepts job but never runs). Token via \texttt{--token-file} (JSON), never CLI/env in repo.

\section{Hardware Measurement Campaign}

The campaign protocol, the raw measurement files and the recomputation scripts are
version-controlled \cite{QOI-hardware}; every number quoted below is regenerated from those
files by \texttt{scripts/summarize\_hardware\_run.py}, which also runs in CI.

\subsection{Platform Constraints}

IBM Quantum open plan (three separate instances --- separate instance identifiers, separate 600 s quotas):
\begin{itemize}
\item No Runtime Sessions: HTTP 400 code 1352 ("not authorized"). The loop script requests a session and falls back to job mode; the measured queue wait proves the fallback (the summary records \texttt{session\_requested=true} but \texttt{execution\_mode=job}).
\item Job-mode queue wait: 6.8--137.5 s (avg 17.9 s/iter).
\item QPU time: 2.96 s/iter (3 circuits + anchor $\times$ 1024 shots).
\item 1000 iterations $\approx$ 2960 s QPU $>$ 600 s quota $\Rightarrow$ open plan infeasible for full protocol.
\end{itemize}

\subsection{Run 1: 68 Iterations (\texttt{ibm\_fez}, seed 20261027)}

\textbf{Scope note.} The QPU-side quantities in this run (Bell parity, $c$ drift, per-channel
decoding error, and the decoded $(A, f, \varphi)$ values themselves) are hardware facts and are
unaffected by the fluid implementation. The \emph{trajectory} is not: the mode-argmax sequence
is a property of the fluid, and the fluid that generated it was the Python port described in
\S\ref{sec:correction}, whose kinetic step was not a Fourier transform. Two further defects
affect the trajectory only. The fed-back frequency was on the wrong scale
($f_{\rm dec}/0.5$ instead of $f_{\rm dec}/0.5 \times 15$, the mode index living in
$[0,15]$), and the simulator's amplitude channel decoded $P(q_0=1)$ where the bridge decodes
$P(q_0=0)$, which inverts the amplitude. Both are corrected in the current code. The
histogram is retained for traceability and is superseded by the corrected deterministic
simulation (\S\ref{sec:quantum-noise}); a hardware re-run with the corrected fluid requires a
plan instance with sufficient quota.

\begin{table}[h]
\centering
\caption{68-iteration loop summary.}
\label{tab:loop68}
\begin{tabular}{lcc}
\toprule
Metric & Value \\
\midrule
Total QPU time & 201 s \\
Total wall time & 1214.8 s \\
Avg queue wait/iter & 17.9 s \\
Avg QPU time/iter & 2.96 s \\
$c$ range & 0.9170--0.9531 \\
$c$ mean & 0.9350 \\
$c$ drift & +0.0361 \\
Mode diversity & 7 modes (\#0: 27, \#7: 10, \#6: 9, \#15: 9, \#14: 6, \#1: 4, \#2: 3) \\
Feedback decode & the $c$-corrected decode (the one \S\ref{sec:scope} shows to be the wrong one) \\
$f$ feedback scale & $f_{\rm dec}/0.5$ --- the mode index lives in $[0,15]$, so the fed-back value was compressed by $15\times$ \\
Classical control & not recorded for this run \\
\bottomrule
\end{tabular}
\end{table}

Per-channel decoding (mean absolute error over 68 iterations):

\begin{table}[h]
\centering
\caption{Decoding accuracy: raw vs $c$-corrected.}
\label{tab:decoding68}
\begin{tabular}{lccc}
\toprule
Channel & Raw MAE & $c$-corrected MAE & Note \\
\midrule
$A$ & 0.026 & 0.212 & \textbf{Correction harms} (asymmetric noise) \\
$f$ & 0.019 & 0.013 & Improvement (34\%) \\
$\varphi$ & 0.430 & 0.261 & Improvement (39\%), bimodal phase (0/$\pi$) \\
\bottomrule
\end{tabular}
\end{table}

Deterministic proof: for ideal measurement ($A_\text{true}=0.3$), $c$-correction yields error $-0.300$ (saturation). The $c$-indicator tracks drift; it is \emph{not} a decoding correction.

\subsection{Run 2: 16-Mode Argmax (6 trials, 1024 shots)}

\begin{table}[h]
\centering
\caption{16-mode argmax on hardware.}
\label{tab:16mode}
\begin{tabular}{lcc}
\toprule
Metric & Quantum & Classical Control \\
\midrule
Bell parity & 93.95\% & --- \\
Argmax accuracy & 0/6 (Wilson 95\% CI [0, 0.39]) & 6/6 (100\%) \\
Quantum winners & Creativity 4$\times$, Logic 2$\times$ & Brain 0$\times$ \\
\bottomrule
\end{tabular}
\end{table}

The 1.85\% gap between Brain (0.096210) and Creativity (0.094434) requires $\sim$492k shots for 3-$\sigma$ statistical separation --- 481$\times$ current 1024. On open plan: $492\text{k shots} \times 0.74\text{ s/circuit} \approx 364\text{ s QPU}$ (single circuit time). \textbf{Fundamental limit:} open plan cannot resolve the mode competition quantum-mechanically; the shot-noise floor (0.137$\sigma$) is the primary barrier, not the quota.

\subsection{Run 3: 20 Iterations Raw Feedback (seed 20261026)}

No $c$-correction, raw decoded values fed back.

\begin{table}[h]
\centering
\caption{20-iteration raw feedback vs classical control.}
\label{tab:raw20}
\begin{tabular}{lcc}
\toprule
Metric & Quantum (raw) & Classical \\
\midrule
Modes visited & 5 (\#0, \#6, \#7, \#8, \#15) & 1 dominant (949/1000) + 5 others \\
Mean $A$ MAE & 0.036 & --- \\
$c$-corrected $A$ MAE & 0.176 & --- \\
\bottomrule
\end{tabular}
\end{table}

The classical control column is the CPU reference of \S\ref{sec:classical-control}: with the
corrected fluid it does not lock to a single mode (949 of 1000 iterations in the dominant
mode, 5 other modes visited), so the fluid has intrinsic mode-switching dynamics on its
own. The quantum branch does not create that diversity \emph{ex nihilo}, and it does not
explain it either --- see \S\ref{sec:control} for the control comparison that settles this.

\subsection{Run 4: Second 20-Iteration Run (seed 20261027)}

\begin{itemize}
\item \texttt{ibm\_kingston}: 52 jobs queued, 900 s timeout $\to$ \texttt{RuntimeJobTimeoutError}, job cancelled cleanly (timeout handler validated). This is a platform observation on a different device, not a replication.
\item \texttt{ibm\_fez}, second 20-iteration run (seed 20261027): 3 modes (\#0, \#6, \#7), 60 s QPU time, 393.0 s wall time, avg wait 19.6 s/iter.
\end{itemize}

\subsection{CPU Control Comparison: Does the Quantum Channel Cause the Mode Switching?}
\label{sec:control}

The runs above cannot answer this on their own: each is a single trajectory, and the fluid
that produced them is no longer the fluid in the repository. The control comparison
(\texttt{scripts/control\_comparison.py}, raw output in \texttt{measurement\_cpu/}) runs the
same loop under four kinds of feedback and compares the mode statistics against the
classical control. Per branch: 200 noise seeds at a fixed initial state, plus 30 varying
initial states at 6 noise seeds each, at both hardware run lengths (20 and 68 iterations),
with both the corrected solver and the pre-fix one.

The classical control does not lie inside the noise band. With the corrected solver at 20
iterations, the noise-free control visits 3 modes with a Brain proportion of 0.35 and a
marginal entropy $H/H_{\max} = 0.395$, while every noisy branch concentrates on a single
mode (Brain 0.85--0.95, $H/H_{\max}$ 0.07--0.19). Measurement noise does not create
structure here: it collapses the distribution onto the dominant mode.

The decisive test replays the hardware's \emph{own recorded feedback} through the same loop.
With the frequency scale the hardware actually used, that feedback sits outside the noise
band on all four statistics. With the corrected frequency scale, the same recorded feedback
lands inside the noise band on all four statistics, for all three hardware runs.

\begin{table}[h]
\centering
\caption{Hardware feedback replayed through the pre-fix solver, 20 iterations. The noise band is the 1024-shot branch (median [5--95\%]).}
\label{tab:control}
\begin{tabular}{lcccc}
\toprule
Feedback interpretation & Modes & Brain & $H/H_{\max}$ & Transient $H$ \\
\midrule
Noise band (1024 shots) & 1 [1,1] & 1.00 [1.00,1.00] & 0.00 [0.00,0.00] & 0.000 [0.000,0.000] \\
Classical control & 1 & 1.00 & 0.00 & 0.000 \\
Hardware feedback, as the hardware fed it & 5 & 0.35 & 0.516 & 0.090 \\
Hardware feedback, corrected $f$ scale & 1 & 1.00 & 0.00 & 0.000 \\
\bottomrule
\end{tabular}
\end{table}

The multi-mode pattern of the 2026-10-27 runs was therefore produced by the frequency-scale
defect, not by the quantum channel. With the scale corrected, the hardware's own
measurements are indistinguishable from the classical control plus measurement noise. The
permutation branch supports the same reading from the other side: reordering the recorded
feedback changes the statistics far less than rescaling it does. For the 20-iteration run
the recorded order sits near the middle of the permutation distribution (mode count at the
60th percentile, marginal entropy at the 47th, transition entropy at the 84th), whereas
against the noise band the same measurement is at the 100th percentile on modes and on
entropy and at the 0th on the Brain proportion --- outside the band on all four. The
\emph{values} of the recorded feedback carry the pattern; the order in which the hardware
happened to produce them does not.

\subsection{Scope of the Claims}

\label{sec:scope}

This is a measurement paper. The claims it supports are: (i) a quantum-classical feedback loop
runs on a real superconducting QPU under an open plan, with every iteration's raw counts,
job IDs and timestamps version-controlled; (ii) the symmetric depolarisation model
$p = c\,p + (1-c)/2$ does \emph{not} describe the observed readout error, and correcting with
the live Bell parity makes the amplitude channel worse; (iii) the mode-competition question
is below the shot-noise floor of the available hardware; (iv) the mode histograms of the
2026-10-27 runs are not evidence for the quantum channel, because with the frequency scale
corrected the same recorded measurements reproduce the classical control's behaviour and lie
inside the measurement-noise band (\S\ref{sec:control}).

The claims it does \emph{not} support: that the system exhibits emergent intelligence, that the
16 modes are cognitive in any validated sense, that the quantum channel outperforms the
classical one, or that the quantum feedback causes the mode switching. The mode names (Brain,
Creativity, Logic, \dots) are labels of an architectural observable space inherited from the
original design; they are not claims about cognition. The memory and reconsolidation
architecture (\S\ref{sec:qam}) is an explicit design proposal and is marked \emph{not
validated} wherever it appears. The 16-mode argmax measurement (6 trials, 0/6 agreement) is
reported rather than interpreted: at $N=6$ it cannot separate "no quantum effect" from "a
small quantum effect".

The bridge's encoding limits are stated in \S\ref{sec:bridge-limits}: the frequency channel is
unambiguous only on $[0,0.5]$, the two-qubit form carries no information beyond the
single-qubit form, and the three channels reduce to one and the same measurement --- so the
"phase" is a rotation angle, not a quantum phase.

\section{A Correction to the Fluid Port}

\label{sec:correction}

An independent review of this repository found that the Python fluid port
(\texttt{scripts/quantum\_loop\_sim.py}, \texttt{scripts/gpe\_fft\_measurement.py}) applied the
kinetic phase multiplier \texttt{KPH} in index space and then applied an inverse transform, without
the forward transform that defines the operator. The operator was therefore not the Gross-Pitaevskii
kinetic step, and the four-branch simulation as well as the hardware loop's fluid ran on it.

The Rust crate \texttt{crates/fluid-stratos} was never affected: it performs
$\text{fft} \to \text{phase} \to \text{ifft}$ with the $1/n^2$ scaling. The corrected Python port
now matches the Rust solver on the same initial state to $1.1\times 10^{-15}$ in
$\max|\Delta\rho|$ (parity test, \S\ref{sec:parity}).

The correction changes the numbers in the direction of \emph{stronger} structure, not weaker:
the density redistribution grows from $2.766$ to $2.780$, the 16-mode energy motion from $8.444$
to $9.420$, the quantum branch's marginal concentration from $H/H_{\max}=0.4355$ to $0.5268$, and
the total-variation distance between the quantum and classical branches from $0.2960$ to $0.3730$.
What it invalidates is one specific inference: that the classical control collapses to a single
mode. With the corrected operator the classical branch visits 949/1000 in one mode and 5 others,
so mode diversity alone does not separate the quantum channel from the fluid's own dynamics.
The separating quantity is therefore the total-variation distance and the marginal concentration,
not the mode count (\S\ref{sec:quantum-noise}).

Two further corrections, applied after the numbers above were measured, reduce that separation
substantially and are the reason the figures in \S\ref{sec:quantum-noise} and
\S\ref{sec:classical-control} differ from the ones quoted here. The amplitude channel decoded
$P(q_0=1)$ where the hardware bridge decodes $P(q_0=0)$, which inverts the amplitude; with the
inverted channel the quantum branch's concentration was $H/H_{\max}=0.5268$ and the
total-variation distance $0.3730$. With the amplitude decoded as the bridge decodes it, the
quantum branch's concentration is $0.2198$ against the classical branch's $0.0993$, and the
total-variation distance is $0.0970$. The frequency normalisation defect of
\S\ref{sec:bridge-limits} was fixed at the same time. The corrected figures are the ones the
repository now regenerates and the ones this paper reports.

\section{Parity Tests}

\label{sec:parity}

Two parity tests guard the two implementation seams that the correction above exposed. Both run
in CI on every commit.

\paragraph{Rust versus numpy.} The same deterministic initial state is advanced for 200 steps by
the Rust solver and by the corrected numpy port; the maximum density difference must stay below
$10^{-12}$ (measured: $1.105\times 10^{-15}$). Without this test the port can silently diverge
again, because the two implementations look interchangeable at the call site.

\paragraph{Serial versus parallel.} The parallel solver must produce the same density as the
serial one after 200 steps within $10^{-12}$. The parallel path used to return silently when the
field shapes did not match, which is exactly the staleness condition the CHANGELOG claims to have
fixed; it now fails loudly instead.

\section{QDS Validation (Quantum Delayed Synchronization)}

The QDS measurement is carried out in its own repository \cite{QDS}; the numbers below are
taken from that repository's \texttt{results/multirun\_stats.json}.

Hahn-echo $T_2$ sweep on \texttt{ibm\_kingston} qubit 91, 5 independent runs,
1024 shots per delay, $\tau \in \{0, 10, 25, 50, 100, 200, 400, 800\}\ \mu$s.
Source of truth: \texttt{quantum-delayed-sync/results/multirun\_stats.json}.

\begin{table}[h]
\centering
\caption{QDS effective $T_2$ per run (qubit 91, \texttt{ibm\_kingston}).}
\label{tab:qds}
\begin{tabular}{lll}
\toprule
Run & Job ID & $T_{2,\text{eff}}\ (\mu\text{s})$ \\
\midrule
1 & \texttt{dba0q2gc77vs738252kg} & 77.53 \\
2 & \texttt{dba0u8hgv2cc73e66p10} & 43.08 \\
3 & \texttt{dba0v8hgv2cc73e66q5g} & 140.75 \\
4 & \texttt{dba1000c77vs73825a6g} & 117.75 \\
5 & \texttt{dba10os24mrc7384npq0} & 69.06 \\
\midrule
Mean $\pm$ SD & --- & $89.63 \pm 39.18$ \\
SEM & --- & $17.52$ \\
\midrule
Calibrated $T_1$ / $T_2$ (IBM, at measurement time) & --- & $254.16$ / $121.26$ \\
Run 1 fit, stderr & \texttt{dba0q2gc77vs738252kg} & $77.53 \pm 12.89$ \\
\bottomrule
\end{tabular}
\end{table}

The measured echo-limited $T_2$ (77.5 $\mu$s for the evaluated run, 89.6 $\mu$s
multi-run mean) sits below IBM's calibrated 121.3 $\mu$s. The QDS repository
documents this deficit as a systematic effect (a wrong $dt$ assumption in the
calibration), not a measurement error; the $\pm 39.2\ \mu$s spread across the five
runs is the run-to-run dispersion of that estimate, not fit uncertainty
(run 1: $\pm 12.89\ \mu$s). The coherence window therefore supports the anchor
protocol at $< 100\ \mu$s, but the dispersion must be carried as an error budget
rather than a single number.

\subsection{Quantum Anchor Protocol}

The anchoring protocol is specified in its own repository \cite{QuantumAnchor}.

Post-pulse coherent-state anchoring on IBM/IQM:
\begin{enumerate}
\item Prepare target state via calibrated pulse.
\item Apply echo sequence at variable delay $\tau$.
\item Measure overlap with reference; fit exponential decay.
\item Anchor accepted if fidelity $> 0.95$ at protocol $\tau$.
\end{enumerate}
Validated on \texttt{ibm\_fez} and IQM denominators; enables QAM state initialization.

\subsection{Wave-Quantum Encoding}

The frequency-to-correlation encoding and its X-basis no-go result live in a separate
repository \cite{WaveQuantum}; we do not reproduce that measurement here.

Frequency-to-correlation mapping: encode $\omega$ in RY angle, measure ZZ correlation in X-basis. Pearson $r = 0.87$--$0.999$ across frequencies. X-basis ZZ theorem proves no-go for direct frequency discrimination without tomography.

\subsection{QAM: Quantum Anchored Memory (\emph{not validated in this work})}
\label{sec:qam}

The QAM architecture is a design proposal assembled in a separate repository \cite{QAM}.
\label{sec:qam}

QAM fuses three layers:
\begin{enumerate}
\item \textbf{Wave dynamics}: $\psi(t) = A e^{-\gamma t} \cos(2\pi f t + \varphi)$ with hardware-measured $(A,f,\varphi,\gamma)$.
\item \textbf{Quantum Anchor}: coherent-state initialization + periodic re-anchoring (fidelity threshold).
\item \textbf{Microscope Memory} \cite{MicroscopeMemory}: 9-level zoom (D0--D8), a layered architecture, binary mmap, SIMD, spine bus REST/MCP, Merkle integrity. Its native Rust cognitive gate is not part of this repository and is not built or tested here.
\end{enumerate}

The spine bridge exposes \texttt{POST /recall}, \texttt{POST /store}, \texttt{GET /stats} with zero JSON overhead (binary frames). Its storage statistics are a property of that separate installation, not a result of this work, and are therefore not tabulated as measurements here.

\subsection{scs-quantum: TKR Encoding}

The time-knowledge representation is implemented in a separate repository
\cite{scsQuantum}; its Bell reference value is quoted for context and is not reproduced here.

Time-Knowledge Representation (TKR): frequency-domain wave-packet encoding on superconducting QPUs. Bell reference 99.3\%, GHZ ideal. Provides the encoding foundation for QAM wavepackets.

\section{Results Synthesis}

\begin{table}[h]
\centering
\caption{Unified results across all modules.}
\label{tab:unified}
\begin{tabular}{llcc}
\toprule
Module & Key Metric & Value & Status \\
\midrule
Fluid STRATOS & Density redistribution $\max|\Delta\rho|$ (500 steps) & $2.780$ & Measured (30 tests) \\
Fluid STRATOS & Coherence (500 steps) & $0.452$ & Measured \\
Quantum Bridge & Bell parity (hardware) & $93.95\%$ & Measured \\
Quantum Bridge & 16-mode argmax accuracy & 0/6 (Wilson 95\% CI [0, 0.39]) & Measured (6 trials) \\
Quantum Bridge & $A$ raw MAE & $0.026$ & Measured \\
Quantum Bridge & $c$-corrected $A$ MAE & $0.212$ & \textbf{Worse} \\
Quantum Bridge & $f$ raw MAE & $0.019$ & Measured \\
Quantum Bridge & $c$-corrected $f$ MAE & $0.013$ & Improvement \\
Quantum Bridge & $\varphi$ raw MAE & $0.430$ & Measured \\
Quantum Bridge & $c$-corrected $\varphi$ MAE & $0.261$ & Improvement \\
Hardware Loop & 68-iter mode diversity & 7 modes & Measured (trajectory superseded, \S\ref{sec:control}) \\
Hardware Loop & Classical control modes & 1 dominant + 5 others & CPU reference (\S\ref{sec:classical-control}) \\
Hardware Loop & Recorded feedback vs measurement-noise band & inside on 4/4 statistics & Measured (\S\ref{sec:control}) \\
QDS & $T_2$ (qubit 91) & $89.63 \pm 39.18\ \mu\text{s}$ & 5 runs \\
Quantum Anchor & Fidelity threshold & $> 0.95$ & Protocol defined \\
Wave-Quantum & Pearson $r$ (freq $\to$ corr) & $0.87$--$0.999$ & Measured \\
\bottomrule
\end{tabular}
\end{table}

The table above contains only quantities measured or recomputed in this work. The architectural
modules that are \emph{not} validated here are listed separately, with their status stated
explicitly:

\begin{table}[h]
\centering
\caption{Architectural modules referenced by the framework, not validated in this work. Their measurements live in their own repositories and are not reproducible from this one.}
\label{tab:notvalidated}
\begin{tabular}{lll}
\toprule
Module & Claimed property & Status here \\
\midrule
Quantum Anchor & Coherent-state anchoring protocol & Protocol only \\
Wave-Quantum & Frequency-to-correlation encoding & Not reproduced \\
QAM & $\psi(t)$ + memory + anchor integration & Design proposal \\
Microscope Memory & Gate tests, archive size & Not reproduced \\
scs-quantum & TKR wave-packet encoding & Not reproduced \\
\bottomrule
\end{tabular}
\end{table}

\section{Discussion}

\subsection{Asymmetric Readout Noise}
The measurement floor ($P_\text{min} \approx 0.037$) and ceiling ($P_\text{max} \approx 0.93$) create asymmetric error: $p_\text{meas} = c p + (1-c)/2$ assumes symmetric depolarization, which fails. The $c$-indicator remains useful for \emph{monitoring} drift (0.917--0.953), but \emph{correction amplifies error} for amplitude.

\subsection{Bimodal Phase Dynamics}
\label{sec:phase}
$\varphi$ clusters near $0$ and $\pi$ (boundary of dynamical domain). At boundaries, readout noise $\pm \delta P$ maps to large $\delta\varphi$ via $\arcsin$ nonlinearity. Phase unwrapping or alternative encoding (e.g., $\sin\varphi, \cos\varphi$ via 2-qubit tomography) needed.

\subsection{16-Mode Discrimination Gap}
The Brain vs. Creativity energy gap (1.85\%) requires $\sim$492k shots for 3-$\sigma$ statistical separation --- 481$\times$ current 1024. On open plan: $492\text{k shots} \times 0.74\text{ s/circuit} \approx 364\text{ s QPU}$ (single circuit time). \textbf{Fundamental limit:} open plan cannot resolve the mode competition quantum-mechanically; the shot-noise floor (0.137$\sigma$) is the primary barrier, not the quota.

\subsection{Limits of the Bridge Encoding}
\label{sec:bridge-limits}

Three properties of the encoding constrain what the bridge can measure, and they are
properties of the protocol rather than of the device.

\emph{The frequency channel is unambiguous only on $[0,0.5]$.} The two-qubit interference
decode is $f = \arcsin(\sqrt{P})/\pi$ and $\sin^2(\pi f)$ is periodic, so a frequency above
$0.5$ folds back. The fluid's $f$ is an energy-weighted average of the 16 mode indices and
therefore lives in $[0,15]$; it must be normalised into $[0,0.5]$ before encoding and scaled
back after decoding. On the recorded hardware files the encoded value was 1.24--3.81, and
111 of 111 iterations folded back. This is a protocol defect, not a hardware limit, and it is
fixed in the current code.

\emph{Two qubits buy no advantage here.} The frequency and phase circuits are an \texttt{RY} on
the control qubit plus a \texttt{CNOT} that copies it into the readout qubit. The probability
law is $P(q_0=1) = \sin^2(\theta/2)$, the same law as a single-qubit rotation. The second
qubit adds circuit size, and with it QPU time and error accumulation, but no information.

\emph{$A$, $f$ and $\varphi$ are the same measurement.} All three reduce to one rotation angle
$\theta$, a measurement of $P(q_0=1)$, and the inverse of $\sin^2(\theta/2)$. $A$ and
$\varphi$ are literally the same inverse, $2\arcsin(\sqrt{P})$; $f$ is that form scaled by
$1/\pi$. The phase is therefore not a quantum phase --- it is a rotation angle measured with
the same statistics as the amplitude, and it is the channel whose fluid-side signal is
bimodal (\S\ref{sec:phase}).

\subsection{Quantum Noise as Perturbation, Not Structure}
\label{sec:quantum-noise}
Simulation with the corrected fluid (numpy FFT, 4 branches, 1000 iterations, deterministic seed):
quantum-loop $H/H_{\max} = 0.2198$ (top mode 87.4\% over the last 500), classical control
$H/H_{\max} = 0.0993$ (949/1000 in one mode, 5 other modes visited), uniform null model 0.9962,
dead GPE 0.0000. The total-variation distance between the quantum and the classical branch is
0.0970. These figures are regenerated by \texttt{scripts/quantum\_loop\_sim.py} and committed as
\texttt{measurement\_cpu/quantum-loop-results.json}.

The two dominant-mode proportions differ by $\hat p = 0.949$ vs $0.874$; the standard error of
that difference is $\sigma_{\Delta} = 0.0126$, i.e. a separation of roughly $6\sigma$ over
1000 iterations. The quantum channel measurably changes the marginal mode distribution, and the
separation is far smaller than the $22\sigma$ and TV distance of $0.3730$ that were reported
before the amplitude channel was corrected. What the data do \emph{not} show is
that the quantum path identifies the correct mode: the 16-mode argmax test fails at 0/6.
And what the control comparison of \S\ref{sec:control} shows is stronger still: the mode
switching of the hardware runs is not produced by the quantum channel at all.
So the quantum channel is a structured perturbation of the fluid's own mode dynamics, not a source
of modal information.

\subsection{Biological Plausibility of Reconsolidation}
QAM's $\psi(t)$ decay + anchor reconsolidation mirrors memory reconsolidation windows \cite{Nader2000}. The 9-level zoom (D0--D8) maps to hippocampal-cortical hierarchy. Cross-validation invited.

\subsection{Open Plan vs. Dedicated Session}
Open plan: no sessions, job queue measured, quota hard limit. Dedicated/flex plan: session reservation, lower latency, higher quota. \textbf{1000-iteration protocol requires dedicated plan.}

\section{Reproducibility \& Artifacts}

\begin{itemize}
\item \textbf{Repository}: \url{https://github.com/silentnoisehun/quantum-organic-intelligence} (public)
\item \textbf{Raw data}: \texttt{measurement\_raw/} (104 JSON files, job IDs, timestamps, version-controlled)
\item \textbf{Reports}: \texttt{docs/MEASUREMENTS-HARDWARE-20261027.txt}, \texttt{RUNLOG-*.txt}
\item \textbf{Summary script}: \texttt{scripts/summarize\_hardware\_run.py} (runs in CI)
\item \textbf{Rust crates}: \texttt{fluid-stratos} (12 tests), \texttt{quantum-bridge} (18 tests, \texttt{bridge-cli} binary)
\item \textbf{Python env}: \texttt{qoi-venv} (qiskit 2.6.1, qiskit-ibm-runtime 0.51.0, numpy 2.5.4)
\item \textbf{Commands}:
\begin{verbatim}
# Local validation (no token)
python scripts/hardware_measure.py --local --shots 256 --mode-trials 2

# Hardware (requires IBM_QUANTUM_API_TOKEN)
python scripts/hardware_measure.py --backend ibm_fez --shots 1024 --mode-trials 6
python scripts/quantum_loop_hardware.py --backend ibm_fez --iterations 100 --shots 1024 --anchor-every 5 --token-file <token-file>.json
\end{verbatim}
\item \textbf{CI}: 7 consecutive green runs (fmt, clippy, tests, measurement scripts, paper PDF)
\end{itemize}

\section{Conclusion}

We built and measured a quantum-classical feedback system on IBM superconducting hardware and audited it end to end. The loop runs (111 documented iterations), but the quantum path does not outperform the classical reference for mode discrimination (0/6, Wilson 95\% CI [0, 0.39] against a 949/1000 dominant classical reference). The $c$-indicator tracks hardware drift but fails as correction for amplitude due to asymmetric readout noise. Phase bimodality dominates $\varphi$ error. The 16-mode gap is statistically unresolved; the shot-noise floor (0.137$\sigma$) is the primary barrier, not the quota.

The audit also overturned the framework's own central inference. The mode switching of the
2026-10-27 runs was attributed to the quantum feedback; the control comparison shows that it
was produced by a protocol defect in the frequency scale, and that the hardware's own
recorded measurements, replayed correctly, are indistinguishable from the classical control
plus measurement noise. What the quantum channel measurably does is perturb the marginal
mode distribution (a $6\sigma$ shift in the dominant-mode proportion, TV distance 0.0970 at
1000 iterations) --- it does not identify a mode, and it does not cause the switching.

Complementary modules (QDS, Quantum Anchor, Wave-Quantum, QAM, Microscope Memory, scs-quantum) provide independently validated building blocks for a neuromorphic quantum memory architecture. The full stack is open: code, data, measurements, and limits.

\textbf{Future work:} a dedicated IBM plan for 1000 iterations, to be run only if the
pre-registered condition in \texttt{docs/QPU-RUN-DECISION-RULE.txt} fires; an asymmetric noise
model for mitigation; phase-unwrapped encoding; QAM biological cross-validation; IQM denominator comparison.

\section*{Acknowledgments}

We thank the IBM Quantum team for open-plan access. This work was self-funded; sponsor: \url{github.com/sponsors/silentnoisehun}.

\bibliographystyle{ACM-Reference-Format}
\bibliography{references}

\appendix

\section{Appendix A: Repository Index}

\begin{table}[h]
\centering
\caption{All 7 repositories in the QOI ecosystem.}
\label{tab:repos}
\begin{tabular}{lll}
\toprule
Repository & Focus & Key Result \\
\midrule
quantum-organic-intelligence & GPE + Bridge + Loop & 111 hardware iterations \\
quantum-anchor & Coherent-state anchoring & Protocol (not reproduced here) \\
wave-quantum & Freq-to-correlation encoding & Not reproduced here \\
quantum-delayed-sync & Hahn-echo $T_2$ sweep & $89.63 \pm 39.18\ \mu\text{s}$ (5 runs) \\
quantum-anchored-memory & $\psi(t)$ + Memory + Anchor & Design proposal \\
scs-quantum & TKR wave-packet encoding & Not reproduced here \\
microscope-memory & 9-level cognitive memory & Not reproduced here \\
\bottomrule
\end{tabular}
\end{table}

\section{Appendix B: Hardware Loop Parameters}

\begin{verbatim}
Backend: ibm_fez (156 qubit, active)
Shots: 1024 per circuit
Circuits/iter: 3 (A, f, phi) + Bell anchor every 5
Anchor update: live c fed to bridge-cli decoder
Timeout: 900 s per job (clean cancel on expiry)
Quota guard: pre-submission service.usage() check
Token: --token-file (JSON), never in repo/env
Venv: qoi-venv (qiskit>=2.0, qiskit-ibm-runtime>=0.50)
\end{verbatim}

\section{Appendix C: Classical Control Baseline}
\label{sec:classical-control}

With the corrected operator the classical branch is not a single-mode attractor: over 1000
iterations it stays in one mode in 949 cases and visits 5 other modes. The fluid therefore
has intrinsic mode-switching dynamics on its own, and the mode count alone cannot attribute
diversity to the quantum channel.

The quantity that separates the two branches is the distribution, not the mode count. The
dominant-mode proportion is 0.874 (quantum) versus 0.949 (classical); the standard error of
the difference is 0.0126, a separation of roughly $6\sigma$. The total-variation distance
is 0.0970. The uniform null model ($H/H_{\max}$ 0.9962) and the dead-GPE model (0.0000)
bracket the result, with the quantum branch (0.2198) much closer to the classical fluid
(0.0993) than to the uniform distribution. The channel is a perturbation of the classical
fluid, not a source of information --- and, as \S\ref{sec:control} shows, not the cause of
the mode switching either.

These figures are regenerated by \texttt{scripts/quantum\_loop\_sim.py} and committed as
\texttt{measurement\_cpu/quantum-loop-results.json}; CI regenerates the file and compares it
against the committed version. The comparison requires the mode sequences and the histograms to
be exactly equal and allows a relative tolerance of $10^{-9}$ on the floating point values,
because the coherence floats differ by about $10^{-13}$ between platforms. The mode sequence
itself is not close to a tie: the smallest decision margin (top one minus top two mode energy)
over the 1000 iterations is $1.4\times10^{-4}$. They supersede an earlier set of values (dominant proportions
0.607 and 0.968, $22\sigma$, TV 0.3730) that were computed with an amplitude channel which
decoded $P(q_0=1)$ where the hardware bridge decodes $P(q_0=0)$, and therefore fed the
amplitude back inverted.

\end{document}